\section{Discussion}
\label{sec:discussion}
\paragraph{When does modality-local security earn its cost?}
The most favorable setting is a legacy-compatible interaction path that
already carries audio to the intended application but lacks a shared secure
agent API or trustworthy media intermediary. Audio-carried control lets the
end applications establish a protected exchange through that available path.
Its cost is additional endpoint deployment, airtime, scheduling, and codec
sensitivity. A business that controls both endpoints of a suitable authenticated
API may obtain the same authorization properties more simply without a modem.
The evaluation must identify this boundary rather than assume universality.

\paragraph{Multimodal and physical AI are a motivation, not evaluated scope.}
A robot observing a sign, an agent hearing a remote speaker, and a model
reading rendered text all face questions about whether perceptual content
may authorize effects. A signed optical symbol, an authenticated acoustic
control exchange, or another modality-specific mechanism might help connect
an observed interface to an issuer. But copying and live relaying still
separate key possession from physical origin. Extending \sys{} to cameras,
robots, or safety-critical actuators would require additional protocols,
presence assumptions, and experiments. No such implementation is claimed.

\paragraph{Authentication does not imply benevolence.}
A valid agent may request an allowed operation for a harmful purpose, and a
customer may authorize a fraudulent task. Enrollment policies, business
controls, anomaly detection, user confirmation, and liability remain necessary.
The fraud adapter is intentionally able to restrict an authenticated call.
Its interface is implemented, but no evaluated commercial detector or fraud
model is included in these results.

\paragraph{A channel does not certify a conversation.}
The current design protects structured control messages and their endpoint
identity. It does not sign utterances, authenticate transcripts, or bind every
spoken claim to the protected request. Speech summarization must not turn a
model's paraphrase into a new authorization. A business should display the
validated action and result through an appropriate trusted interface and
make human takeover or changed endpoints trigger a new authorization decision.
Secure transfer and multi-party call semantics remain open work.

\paragraph{Identity governance and privacy.}
An authority's certificate quality depends on its registration procedures.
Production adoption needs revocation distribution, short-lived credentials,
issuer selection, recovery, and compromise response. Stable certificates can
also enable correlation across businesses; pairwise identities and minimal
disclosure deserve study. Our prototype is not an identity federation or a
privacy-preserving credential system. Keys and tokens must remain protected
in the agent runtime, and logs and recordings require a deliberate retention
policy even when application payloads are encrypted.

\paragraph{Submission readiness.}
The checked-in prototype demonstrates a working mechanism with explicit
limitations. A full systems-security submission needs the adversarial,
compositional, and real-path evaluation in Section~\ref{sec:evaluation}, plus
independent review of the authentication integration. A networking submission
also needs a strong quantitative account of modem adaptation, scheduling,
congestion/recovery behavior, and deployment tradeoffs. We do not infer
publication-level evidence from a test count or a sub-second local handshake.

\section{Conclusion}
\label{sec:conclusion}
Personal-agent voice interactions expose a gap between plausible conversation
and explicit customer authority. \sys{} demonstrates that existing audio
adapters can carry a mutually authenticated cryptographic control exchange and
that a separate gateway can enforce scoped grants before business execution.
Its QUIC profile recovers deterministic local audio faults without reliable
link framing. The remaining research question is where this binding provides
measurable security and deployment value under realistic paths and adversaries.
The evaluation plan specifies how to test that question without treating
cryptographic endpoint identity as proof of truthful speech or physical origin.
